\documentclass[11pt,a4paper]{article}

\usepackage[utf8]{inputenc}
\usepackage[T1]{fontenc}
\usepackage[english]{babel}
\usepackage{geometry}
\geometry{left=2.5cm,right=2.5cm,top=2.5cm,bottom=2.5cm}
\usepackage{listings}
\usepackage{hyperref}
\usepackage{enumitem}

\title{\textbf{Project Dependencies}}
\author{HandGestures2Bot}
\date{\today}

\lstset{
  basicstyle=\ttfamily\small,
  breaklines=true,
  frame=single,
  numbers=left,
  numberstyle=\tiny,
  showstringspaces=false
}

\begin{document}
\maketitle
\thispagestyle{empty}
\tableofcontents
\newpage

This document provides a list of all tools, libraries, and compilers required to build, run, and develop the Flutter mobile application.

\section{Flutter Mobile Application (\texttt{flutterapp/})}
The mobile app is built using the Flutter framework.

\subsection{Tooling \& Setup}
\begin{itemize}
    \item \textbf{Flutter SDK:} Version \texttt{\^{}3.9.2} or compatible is required.
    \item \textbf{Installation:} From the \texttt{flutterapp/} directory, run the following command to install all Dart and Flutter dependencies:
    \begin{lstlisting}[language=bash]
flutter pub get
    \end{lstlisting}
\end{itemize}

\subsection{Dart/Flutter Packages (\texttt{pubspec.yaml})}
\begin{itemize}
    \item \texttt{file\_picker}: Allows the user to select files from their device.
    \item \texttt{web\_socket\_channel}: Manages the real-time WebSocket connection to the robot for sending commands.
    \item \texttt{camera}: Accesses the device's camera to provide a video stream for gesture processing.
    \item \texttt{permission\_handler}: Requests and manages necessary device permissions (e.g., camera, microphone).
    \item \texttt{shared\_preferences}: Saves persistent key-value data.
    \item \texttt{flutter\_gemma}: Runs the local Gemma AI model for interpreting voice commands. The model file (\texttt{gemma3-1B-it-int4.task}) was downloaded from Kaggle: \url{https://www.kaggle.com/models/google/gemma-3}.
    \item \texttt{speech\_to\_text}: Captures the user's voice and converts it into text for the AI to process.
    \item \texttt{flutter\_tts}: Provides basic, on-device text-to-speech feedback.
    \item \texttt{path\_provider}: Finds correct, platform-specific paths on the device's filesystem for storing data.
    \item \texttt{http}: Makes network requests to the ElevenLabs API for advanced text-to-speech.
    \item \texttt{audioplayers}: Plays the audio files received from the ElevenLabs API.
\end{itemize}

\subsection{Android-Specific Dependencies (\texttt{build.gradle.kts})}
These are installed automatically by Gradle during the build process.
\begin{itemize}
    \item \textbf{Java Development Kit (JDK):} Version 11.
    \item \textbf{Android NDK:} Version specified by the Flutter SDK.
    \item \textbf{MediaPipe:} \texttt{com.google.mediapipe:tasks-vision:0.10.26}. The hand landmarker model was sourced from \url{https://ai.google.dev/edge/mediapipe/solutions/vision/hand_landmarker}.
    \item \textbf{AndroidX Camera:} \texttt{androidx.camera:*}
    \item \textbf{Gson:} \texttt{com.google.code.gson:gson:2.10.1}.
\end{itemize}

\section{Robot-Side Software (\texttt{robot/}, \texttt{ros2\_ws/})}
The robot side handles Bluetooth-based Wi-Fi provisioning, WebSocket communication, and ROS~2-based robot control.

\subsection{Bluetooth Provisioning (Python)}
Bluetooth provisioning enables the robot to receive Wi-Fi credentials via BLE and configure network access.

\subsubsection{Python Packages}
Installed via \texttt{pip3} unless otherwise noted.
\begin{itemize}
	\item \texttt{bluezero}: BLE GATT server library built on BlueZ and GLib.
	\item \texttt{netifaces}: Network interface and IP address utilities.
	\item \texttt{bleak}: BLE client utilities used by helper and diagnostic scripts.
\end{itemize}

\subsubsection{Python Runtime Bindings}
Required for GLib event loop integration.
\begin{itemize}
	\item \texttt{python3-gi}: PyGObject bindings for Python.
	\item \texttt{gir1.2-glib-2.0}: GLib introspection data.
\end{itemize}

\subsubsection{System Services and Tools}
These dependencies must be installed at the OS level.
\begin{itemize}
	\item \textbf{BlueZ}: Bluetooth stack (\texttt{bluetoothd}, \texttt{libbluetooth}).
	\item \textbf{D-Bus}: Required for Bluetooth and network control.
	\item \textbf{NetworkManager}: \texttt{nmcli} used for Wi-Fi configuration.
	\item \textbf{wpa\_supplicant}: Alternative Wi-Fi management backend.
	\item \textbf{dhclient}: DHCP client for IP assignment.
\end{itemize}

\subsubsection{Installation Hint}
The repository provides a helper script that installs core Python dependencies:
\begin{lstlisting}[language=bash]
	pip3 install -U bluezero netifaces
\end{lstlisting}

\subsubsection{Service Scripts}
\begin{itemize}
	\item \texttt{start\_bluetooth\_provisioning.sh}: Launches BLE provisioning.
	\item \texttt{start\_server.sh}: Started automatically after successful provisioning.
\end{itemize}
\newpage
\subsection{ROS~2 Workspace (\texttt{ros2\_ws/})}
The ROS~2 workspace contains the robot control logic and communication server.

\subsubsection{ROS~2 Runtime}
Installed via system package manager or \texttt{rosdep}.
\begin{itemize}
	\item \texttt{rclpy}: ROS~2 Python client library.
\end{itemize}

\subsubsection{ROS~2 Message Packages}
\begin{itemize}
	\item \texttt{geometry\_msgs}
	\item \texttt{std\_msgs}
	\item \texttt{turtlebot3\_msgs}
\end{itemize}

\subsubsection{Python Packages Used by ROS Nodes}
\begin{itemize}
	\item \texttt{websockets}: Asynchronous WebSocket server for app--robot communication.
	\item \texttt{netifaces}: Used for dynamic IP discovery.
	\item \texttt{setuptools}: Required for \texttt{ament\_python} packaging.
\end{itemize}

\subsubsection{Build and Packaging}
\begin{itemize}
	\item \textbf{Build Type:} \texttt{ament\_python}
	\item Dependencies declared in \texttt{package.xml} and \texttt{setup.py}
	\item ROS~2 dependencies resolved via \texttt{apt} or \texttt{rosdep}
\end{itemize}

\subsection{Concrete Robot Provisioning Checklist (Linux)}
\textbf{System Packages:}
\begin{itemize}
	\item BlueZ (Bluetooth daemon and libraries)
	\item D-Bus, GLib, PyGObject
	\item NetworkManager (\texttt{nmcli}) or \texttt{wpa\_supplicant}
	\item ROS~2 base installation (for chosen distribution)
\end{itemize}

\textbf{Python Packages (\texttt{pip3}):}
\begin{itemize}
	\item \texttt{bluezero}
	\item \texttt{netifaces}
	\item \texttt{bleak}
	\item \texttt{websockets}
\end{itemize}

\textbf{ROS~2 Packages (via \texttt{apt}/\texttt{rosdep}):}
\begin{itemize}
	\item \texttt{ros-<distro>-rclpy}
	\item \texttt{ros-<distro>-geometry-msgs}
	\item \texttt{ros-<distro>-std-msgs}
	\item \texttt{ros-<distro>-turtlebot3-msgs}
\end{itemize}


\end{document}
